\section{One coefficient ring for all flat charts}
\label{sec:flat-chart-descent-v053}

\status{FCD1--FCD18 prove the original homogeneous-matrix flat-chart
descent theorem, the finite graded presentation theorem, and the
Noetherian graded-cube model theorem.  The original coefficient ring is
arbitrary.  One coefficient subring, one presentation and one comparison
map serve every chart.  The subsequent construction of this graded
module from actual geometric sections remains separate; no global
\texttt{Perf} or section law is asserted here.  The accepted formal
source is fixed at
\texttt{7cbf7a0c32fca1badf05931879a37c3f8a2b03b4}.}

The section-model argument needs a finite complex that continues to
describe the same sections after every change of coefficients.  The
linear algebra producing that complex was proved in GSL5--GSL10.  Its
remaining input is a suitable model over a Noetherian coefficient ring.
The present section constructs the algebraic part of that model.
The important choice is quite concrete: retain the coefficients of the
presentation \emph{and of certificates for all the chart properties}.
Keeping only the presentation coefficients would give no reason for
flatness to survive over the smaller ring.

All rings are commutative with identity, including the zero ring.
Modules need not be faithful, and no coefficient extension below is
assumed flat.  If $\iota$ is a finite ordered set, write
\[
 S_R=R[X_i\mid i\in\iota],\qquad
 u_i=X_i-1,\qquad X_J=\prod_{i\in J}X_i.
\]
For a matrix $\Phi\in\operatorname{Mat}_{a\times b}(S_R)$,
$\operatorname{coker}\Phi$ means
$S_R^a/\Phi S_R^b$.  Thus the source's
\texttt{QuotSMulTop} $u_i(\operatorname{coker}\Phi)$ is exactly
\[
 Q_i(\Phi)=S_R^a/(\Phi S_R^b+u_iS_R^a).                 \tag{FCD.1}\label{FCD.1}
\]
A $\mathbb Z$-graded module has a direct-sum decomposition; in
particular, every one of its elements has only finitely many nonzero
homogeneous components.  The variables have degree one.  A free
module with basis degrees $d_1,\ldots,d_a$ has degree-$n$ part
$\bigoplus_l(S_R)_{n-d_l}e_l$.

We use the usual elementary properties of tensor products, localization
and projective modules, and Hilbert's basis theorem.  The needed
Artin--Rees, completion and smoothness arguments are given in FCD5.
These conventional algebraic inputs are distinct from
the geometric producers retained after GSL26.  The source map at the
end identifies both the original declarations and the pinned Mathlib
completion results used for this background.

\subsection{Finite certificates attached to one presentation}

% BEGIN FCD1
\paragraph{FCD1: finite relations for a prescribed set of generators.}
\label{claim:FCD1}
Let $B$ be a finitely presented $R$-algebra, and let
$\pi:R[X_1,\ldots,X_m]\twoheadrightarrow B$ be a surjective
$R$-algebra map.  Its kernel is a finitely generated ideal, even when
$R$ is not Noetherian.  Similarly, if $Q$ is a finitely presented
$S$-module and $F\twoheadrightarrow Q$ is a surjection from a
finitely generated $S$-module, its kernel is finitely generated.

\begin{proof}
Choose a finite presentation
$B=R[Y_1,\ldots,Y_t]/(h_1,\ldots,h_r)$.
Lift each image of $Y_l$ to a polynomial $f_l(X)$, and express each
$\pi(X_j)$ by a polynomial $v_j(Y)$ in this presentation.  The
polynomials
\[
 h_k(f(X)),\qquad X_j-v_j(f(X))                         \tag{FCD.2}\label{FCD.2}
\]
belong to $\ker\pi$.  They generate it.  Indeed, if $p(X)$ maps to
zero, then $p(v(Y))$ belongs to $(h_1,\ldots,h_r)$.
After substituting $Y=f(X)$, the polynomial $p(v(f(X)))$ is in
the ideal generated by the first family in \eqref{FCD.2}.
For any polynomial $p$ the difference $p(X)-p(Z)$ lies in
$(X_1-Z_1,\ldots,X_m-Z_m)$: check a monomial by replacing one
factor at a time, and then add.  With $Z=v(f(X))$ this proves the
assertion about $\ker\pi$.

For the module assertion, first let $E=S^a\twoheadrightarrow Q$
be finite free.  Take a known finite presentation
$S^b\to S^c\xrightarrow{q}Q\to0$, and lift the generators in
both directions to maps $\alpha:S^c\to E$ and
$\beta:E\to S^c$.  If $f:E\to Q$ is the given surjection,
these maps satisfy $f\alpha=q$ and $q\beta=f$.
The kernel of $f$ is generated by the columns of $1-\alpha\beta$
and the images under $\alpha$ of the finitely many given relations.
For $z\in\ker f$, write
$z=(1-\alpha\beta)z+\alpha\beta z$; the second term has the
required form because $q\beta z=0$.  Every displayed generator is
in the kernel.  Finally choose a finite free surjection $E\to F$.
The finitely generated kernel of the composite $E\to Q$ maps
surjectively onto $\ker(F\to Q)$.  This proves the general case.
\end{proof}
% END FCD1

% BEGIN FCD2
\paragraph{FCD2: a generalized inverse from the actual projective quotient.}
\label{claim:FCD2}
Let $S\twoheadrightarrow B$ be a surjective algebra map $\pi$,
with kernel $I=(g_1,\ldots,g_r)$.  Suppose $u\in I$ and
$\epsilon:S^a\twoheadrightarrow P$ is $\pi$-semilinear, where
$P$ is projective over $B$.  Suppose also that
\[
 \ker\epsilon=\Phi S^b+uS^a.                          \tag{FCD.3}\label{FCD.3}
\]
There are matrices $\Psi\in\operatorname{Mat}_{b\times a}(S)$
and $E_k\in\operatorname{Mat}_{a\times b}(S)$ such that
\[
 \Phi\Psi\Phi=\Phi+\sum_{k=1}^r g_kE_k.              \tag{FCD.4}\label{FCD.4}
\]

\begin{proof}
Put $p_l=\epsilon(e_l)$.  Semilinearity gives
$\epsilon(y)=\sum_l\pi(y_l)p_l$.  Consequently the map
$F:B^a\to P$, $F(\beta)=\sum_l\beta_lp_l$, is surjective.
Projectivity supplies a $B$-linear section $s:P\to B^a$, so
$Fs=1_P$.  For each $l$, lift $e_l-s(p_l)\in B^a$ to
$y_l\in S^a$.  Its image under $\epsilon$ is zero.  By
\eqref{FCD.3}, write $y_l=\Phi w_l+u w'_l$.  Since $\pi(u)=0$,
the matrix $\Psi$ with columns $w_l$ satisfies
\[
 \overline\Phi\,\overline\Psi=1_{B^a}-sF.           \tag{FCD.5}\label{FCD.5}
\]
Here bars mean application of $\pi$ to the entries.  The columns
of $\Phi$ lie in $\ker\epsilon$, hence $F\overline\Phi=0$.
Multiplying \eqref{FCD.5} on the right by $\overline\Phi$ gives
$\overline\Phi\,\overline\Psi\,\overline\Phi
=\overline\Phi$.  Every entry of
$\Phi\Psi\Phi-\Phi$ therefore lies in $I$.
Express its finitely many entries as sums $\sum_k g_k e_{k,l,j}$
and assemble $E_k=(e_{k,l,j})$.  This is \eqref{FCD.4}.

The section $s$ and its lifted columns were chosen for the given
$P,\epsilon,\Phi$.  No separately chosen quotient module or
identification is introduced.
\end{proof}
% END FCD2

% BEGIN FCD3
\paragraph{FCD3: all certificates on a single chart.}
\label{claim:FCD3}
Fix $i\in\iota$, and suppose the chart data consist of a smooth
$R$-algebra $B$, a surjection $\pi:S_R\to B$ with
$\pi(X_i)=1$, a projective $B$-module $P$, and a surjective
$\pi$-semilinear map $\epsilon:S_R^a\to P$ satisfying
\eqref{FCD.3} with $u=u_i$.
There is a finite family $g_1,\ldots,g_r$ generating $\ker\pi$,
and finite polynomial data $c_k,C_k,C'_k,\Psi,E_k$ such that
\begin{align}
 u_i&=\sum_k c_kg_k,\nonumber\\
 g_k1_a&=\Phi C_k+u_iC'_k,\nonumber\\
 \Phi\Psi\Phi&=\Phi+\sum_k g_kE_k.                   \tag{FCD.6}\label{FCD.6}
\end{align}
Here $C_k$ has size $b\times a$, $C'_k$ has size $a\times a$,
and the last two sizes are as in FCD2.  The quotient $S_R/(g_k)$
is formally smooth over $R$.

\begin{proof}
Smoothness includes finite presentation and formal smoothness.
FCD1 gives the finite generating family for $\ker\pi$.
Since $\pi(u_i)=0$, express $u_i$ in that ideal to obtain $c_k$.
For every $k,l$, semilinearity implies
$\epsilon(g_ke_l)=\pi(g_k)p_l=0$.
Use \eqref{FCD.3} to write
$g_ke_l=\Phi w_{k,l}+u_iw'_{k,l}$; the vectors $w_{k,l}$ and
$w'_{k,l}$ are the columns of $C_k,C'_k$.  This proves the
middle identity of \eqref{FCD.6}.  FCD2 proves the last identity.
Finally the actual map $\pi$ identifies $S_R/(g_k)$ with $B$
as an $R$-algebra, and transports its formal smoothness.  Only
finitely many matrices and entries have been chosen.
\end{proof}
% END FCD3

% BEGIN FCD4
\paragraph{FCD4: formal smoothness expressed by finite polynomial identities.}
\label{claim:FCD4}
Let $S=R[X_i\mid i\in\iota]$ and $I=(g_1,\ldots,g_r)$.
The quotient $A=S/I$ is formally smooth over $R$ if and only if
there are polynomials $a_{ik}$ and $b_{jkl}$ such that, with
\[
 y_i=X_i+\sum_k a_{ik}g_k,
\]
one has
\[
 g_j(y)=\sum_{k,l}b_{jkl}g_kg_l\quad(1\le j\le r).   \tag{FCD.7}\label{FCD.7}
\]
This criterion is valid over an arbitrary coefficient ring.

\begin{proof}
If $A$ is formally smooth, lift $1_A$ across the square-zero
surjection $S/I^2\to S/I$ to an $R$-algebra section
$\sigma:A\to S/I^2$.  Lift each $\sigma(\overline X_i)$ to a
polynomial $y_i\in S$.  Because the reduction of $\sigma$ is the
identity, $y_i-X_i\in I$, giving $a_{ik}$.  Evaluation at $y$
followed by reduction modulo $I^2$ equals the composite
$S\to A\xrightarrow{\sigma}S/I^2$: these maps agree on $R$
and on every variable.  Thus $g_j(y)\in I^2$.
The products $g_kg_l$ generate $I^2$, giving the coefficients
$b_{jkl}$ in \eqref{FCD.7}.

Conversely \eqref{FCD.7} defines an $R$-algebra map
$\sigma:A\to S/I^2$ by $\overline X_i\mapsto\overline y_i$.
Its reduction is the identity.  To verify formal smoothness directly,
let $D\to D/J$ be an $R$-algebra quotient with $J^2=0$, and
let $h:A\to D/J$.  Choose lifts in $D$ of the finitely many
$h(\overline X_i)$.  Evaluation gives $t:S\to D$, with
$t(I)\subseteq J$, hence $t(I^2)=0$.  The induced map
$\overline t:S/I^2\to D$ composed with $\sigma$ lifts $h$.
For a nilpotent ideal $J$, lift successively through
$D/J^{n+1}\to D/J^n$: its kernel $J^n/J^{n+1}$ is square-zero
for $n\ge1$.  This proves the full infinitesimal lifting property.

The coefficients of \eqref{FCD.7} are the finite evidence of formal
smoothness that can be retained in a smaller ring.  A bare assertion
that the original chart was smooth would not by itself provide such
evidence after descent.
\end{proof}
% END FCD4

% BEGIN FCD5
\paragraph{FCD5: the flatness consequence over a Noetherian coefficient ring.}
\label{claim:FCD5}
If $R_0$ is Noetherian, $S_0=R_0[X_i\mid i\in\iota]$, and
$A_0=S_0/I_0$ is formally smooth over $R_0$, then $A_0$ is
flat over $R_0$.

\begin{proof}
Hilbert's basis theorem makes $S_0$ Noetherian.  Let
$\widehat S_0=\varprojlim_{n\ge1}S_0/I_0^n$.
Starting with $1_{A_0}:A_0\to S_0/I_0$, formal smoothness gives
compatible lifts $\sigma_n:A_0\to S_0/I_0^n$.
At the next step the kernel $I_0^n/I_0^{n+1}$ has square zero,
so the lift can be chosen to reduce to the \emph{already chosen}
$\sigma_n$.  The resulting $\sigma:A_0\to\widehat S_0$
is a section of the first projection $p:\widehat S_0\to A_0$.

For completeness, the completion fact used here is that the
$I$-adic completion of a Noetherian ring $S$ is flat over $S$.
Here is its usual module argument.  For an inclusion $N\subseteq M$
of finite $S$-modules, Artin--Rees says that the filtrations
$I^nN$ and $N\cap I^nM$ define the same completion: for a fixed
$c$ and $n\ge c$,
\[
 I^nN\subseteq N\cap I^nM\subseteq I^{n-c}N.
\]
The required inclusion has a short proof.  The Rees ring
$\mathcal R=\bigoplus_{n\ge0}I^nt^n=S[It]$ is a finitely
generated $S$-algebra, hence Noetherian.  The module
$\bigoplus_{n\ge0}I^nMt^n$ is finite over $\mathcal R$,
generated by any finite generating family of $M$ in degree zero.
Its graded submodule
$\bigoplus_{n\ge0}(N\cap I^nM)t^n$ is therefore finite.
Choose finitely many homogeneous generators, all in degrees at
most $c$.  In degree $n\ge c$ they express every element of
$N\cap I^nM$ as a sum of elements of
$I^{n-d}(N\cap I^dM)$ with $d\le c$, each contained in
$I^{n-c}N$.  The other inclusion is immediate.  This proves
the Artin--Rees assertion used here.

The finite-level sequences with left term
$N/(N\cap I^nM)$, middle term $M/I^nM$ and right term
$(M/N)/I^n(M/N)$ are exact.  The transition maps on the left
are surjective.  Consequently taking inverse limits preserves this
short exact sequence.  Explicitly, a compatible right-hand sequence
can be lifted recursively: after any lift at the next level, its
discrepancy from the previously chosen lift is in the left term and
can be corrected using that surjectivity.  Artin--Rees then identifies
the left limit with $\widehat N$.  Thus completion is exact on
finite modules.

A finite presentation $S^b\to S^a\to M\to0$ now identifies
$\widehat M$ with $\widehat S\otimes_SM$: both are the
cokernel of the same matrix $\widehat S^b\to\widehat S^a$.
In particular, for every ideal $J\subseteq S$, which is finite
because $S$ is Noetherian, the map
$\widehat S\otimes_SJ\to\widehat S$ is injective.
The ideal criterion for flatness gives flatness of $\widehat S$.
This argument uses neither a separatedness assumption on $S$
nor an assumption that $I$ lies in its Jacobson radical.

Apply this to $S_0,I_0$.  Since $S_0$ is free over $R_0$,
$\widehat S_0$ is flat over $R_0$ by composition of tensor
functors.  The maps $\sigma,p$ exhibit $A_0$ as an $R_0$-module
retract of $\widehat S_0$.  A retract of a flat module is flat:
after tensoring an injection, its induced map is a direct summand
of an injection.  Hence $A_0$ is $R_0$-flat.
\end{proof}
% END FCD5

% BEGIN FCD6
\paragraph{FCD6: the explicit retract for a dehomogenized cokernel.}
\label{claim:FCD6}
Let $S$ be an $R$-algebra, $I=(g_1,\ldots,g_r)$, and suppose
$S/I$ is flat over $R$.  Suppose $u,\Phi$ and matrices satisfy
the three identities \eqref{FCD.6}, with $u$ in place of $u_i$.
Then
$Q=S^a/(\Phi S^b+uS^a)$ is $R$-flat.  More precisely, $Q$ is
an explicit $R$-module retract of $(S/I)^a$.

\begin{proof}
Write $\operatorname{pr}:S^a\to Q$ for the quotient map and
$q:S^a\to(S/I)^a$ for reduction of the coordinates.
Applying $g_k1_a=\Phi C_k+uC'_k$ to a vector shows that
$g_k$ annihilates $Q$.  Thus $I$ annihilates $Q$, and there is
a well-defined map
\[
 p:(S/I)^a\longrightarrow Q,\qquad p(q(v))=\operatorname{pr}(v).
\]
Define the other map by the concrete formula
\[
 j:Q\longrightarrow(S/I)^a,\qquad
 j(\operatorname{pr}(v))=q(v-\Phi\Psi v).             \tag{FCD.8}\label{FCD.8}
\]
To check that this is well-defined, first apply $1-\Phi\Psi$
to $\Phi w$.  By the last identity of \eqref{FCD.6} the result
is $-\sum_k g_kE_kw$, whose reduction is zero.  Secondly,
$(1-\Phi\Psi)uv$ reduces to zero because
$u=\sum_kc_kg_k\in I$.  These are exactly the two families
of relations in $Q$.

The composite satisfies
\[
 p j(\operatorname{pr}(v))
 =\operatorname{pr}(v-\Phi\Psi v)=\operatorname{pr}(v).
\]
Hence $pj=1_Q$.  A finite direct sum of the $R$-flat module
$S/I$ is flat, and FCD5's retract argument applies.  All maps
were given before any further coefficient change.  Tensoring the
identity $pj=1_Q$ therefore preserves this retract over every
new coefficient ring as well.
\end{proof}
% END FCD6

% BEGIN FCD7
\paragraph{FCD7: descend all finitely many identities to one coefficient ring.}
\label{claim:FCD7}
Let a finite family of polynomials over an arbitrary ring $R$ be given.
There is a Noetherian subring $R_0\subseteq R$ containing their
coefficients.  Their unique coefficientwise lifts to $S_{R_0}$
satisfy every polynomial identity between the original family,
including identities involving substitution and matrix multiplication.
If a specified polynomial is homogeneous of an integer degree $d$,
its lift has that same degree.

\begin{proof}
Let $T\subseteq R$ be the union of the coefficient sets of the
finitely many polynomials, and let $R_0$ be the subring generated
by $T$.  This is the image of a polynomial ring
$\mathbb Z[t_1,\ldots,t_N]$, so it is Noetherian by Hilbert's
basis theorem and passage to a quotient.  The argument also covers
positive characteristic and the zero ring.

Each polynomial lifts by viewing its coefficients as elements of
$R_0$.  The map $S_{R_0}\to S_R$ is injective, since equality
of polynomials means equality of every coefficient.  It commutes
with sums, products and substitution in the variables.  Thus the
difference of the two lifted sides of an identity maps to zero
and is already zero in $S_{R_0}$.  Matrix identities reduce
entrywise to these operations.

For the homogeneity assertion, the coefficient of every monomial
of degree different from $d$ is zero after inclusion in $R$.
Injectivity makes it zero in $R_0$.  This works also for negative
$d$, when the only homogeneous polynomial is zero.  No division
or cancellation in $R$ has occurred.
\end{proof}
% END FCD7

% BEGIN FCD8
\paragraph{FCD8: the original simultaneous flat-chart descent theorem.}
\label{claim:FCD8}
Let $R$ be arbitrary, $\iota$ finite and ordered, and let
$d:\{1,\ldots,a\}\to\mathbb Z$ and
$e:\{1,\ldots,b\}\to\mathbb Z$.  Suppose
$\Phi\in\operatorname{Mat}_{a\times b}(S_R)$ has
$\Phi_{lj}$ homogeneous of degree $e_j-d_l$.
For \emph{every} $i\in\iota$, suppose there exist a smooth
$R$-algebra $B_i$, a surjective $R$-algebra map $\pi_i:S_R\to B_i$
with $\pi_i(X_i)=1$, a projective $B_i$-module $P_i$, and a
surjective $\pi_i$-semilinear map $\epsilon_i:S_R^a\to P_i$
such that
\[
 \epsilon_i(y)=0
 \quad\Longleftrightarrow\quad
 y=\Phi w+(X_i-1)w'\quad\hbox{for some }w,w'.         \tag{FCD.9}\label{FCD.9}
\]
Then there exist one Noetherian subring $R_0\subseteq R$ and
one matrix $\Phi_0$ over $S_{R_0}$ such that its entrywise
image is $\Phi$, its entries have the same degrees $e_j-d_l$,
and
\[
 Q_i(\Phi_0)\ \text{is flat over }R_0
 \quad\text{for every }i\in\iota.                   \tag{FCD.10}\label{FCD.10}
\]
This is the full conclusion of the source theorem
\begin{center}
\texttt{exists\_\allowbreak{}noetherian\_\allowbreak{}subring\_\allowbreak{}}\allowbreak\texttt{homogeneous\_matrix\_flat\_charts}.
\end{center}

\begin{proof}
Choose one set of the supplied data for each of the finitely many
$i$.  Apply FCD3, and then FCD4, to obtain its families
$g_{i,k},c_{i,k},C_{i,k},C'_{i,k},\Psi_i,E_{i,k}$ and the
formal-smoothness polynomials $a_{i,v,k},b_{i,j,k,l}$.
Collect the coefficients of \emph{all} these polynomials and
of all entries of $\Phi$.  The index set, every relation family
and every matrix are finite.  FCD7 therefore gives one $R_0$
and coefficientwise lifts of every family, including $\Phi_0$.
It also preserves the homogeneous degrees of $\Phi$.

Fix any chart $i$.  Its lifted equations \eqref{FCD.7} prove
formal smoothness of
\[
 B_{i,0}=S_{R_0}/(g_{i,1,0},\ldots,g_{i,r_i,0})
\]
by FCD4.  The ring $R_0$ is Noetherian, so FCD5 proves that
$B_{i,0}$ is $R_0$-flat.  All three equations
\eqref{FCD.6} hold over $S_{R_0}$ by FCD7.  Apply FCD6 with
$u=X_i-1$ and this same $\Phi_0$ to obtain \eqref{FCD.10}.
This works for every $i$ without enlarging or changing $R_0$.

Notice the order: the finite chart witnesses are chosen first;
the common coefficient ring is then formed from their union;
all flatness assertions are proved on that one ring.  The input
$R$ itself has never been assumed Noetherian or reduced.
\end{proof}
% END FCD8

\subsection{A finite graded presentation that is correct on every chart}

% BEGIN FCD9
\paragraph{FCD9: three elementary dehomogenization identities.}
\label{claim:FCD9}
Let $M=\bigoplus_{n\in\mathbb Z}M_n$ be graded, and put
$x=X_i$ for a fixed variable.  Then:
\begin{enumerate}
\item multiplication by $x-1$ on $M$ is injective;
\item if a homogeneous $z\in M_n$ satisfies $z=(x-1)w$,
then $x^kz=0$ for some $k\ge0$;
\item for each $y\in M$, and every sufficiently large integer $N$,
there are $y_N\in M_N$ and $v\in M$ with $y=y_N+(x-1)v$.
\end{enumerate}
These assertions do not require multiplication by $x$ to be injective.

\begin{proof}
Write $w=\sum_jw_j$.  The degree-$j$ component of $(x-1)w$
is $xw_{j-1}-w_j$.  If this is zero for every $j$, start below
the finite support of $w$ and use $w_j=xw_{j-1}$ successively;
every component is zero.  This proves the first assertion.

For the second assertion the same recurrence holds except in degree
$n$.  It gives $w_{n-1}=0$, then $z=-w_n$, and
$w_{n+k}=x^kw_n$ for $k\ge0$.  Above the finite support of
$w$ the left side is zero, so some $x^kz$ is zero.

For the third assertion choose $N$ at least every degree in the
support of $y=\sum_jy_j$, and set
\[
 y_N=\sum_j x^{N-j}y_j\in M_N.
\]
The geometric-sum identity
$x^k-1=(x-1)(1+x+\cdots+x^{k-1})$ shows term by term that
$y-y_N$ is a multiple of $x-1$.  The formula also covers $k=0$
by taking an empty sum.  This supplies the required $v$ explicitly.
\end{proof}
% END FCD9

% BEGIN FCD10
\paragraph{FCD10: finite generators up to powers of each chart variable.}
\label{claim:FCD10}
If $M/(X_i-1)M$ is finitely generated over $S_R$, there is a
finite homogeneous family in $M$ such that every $y\in M$ has
some $X_i^ky$ in its $S_R$-span.  If this hypothesis holds for
every $i\in\iota$, one finite homogeneous family works for
all $i$, with the exponent allowed to depend on $i$ and $y$.
More generally, if each chart already has a finite family with
this power-span property, their homogeneous components give one
homogeneous family with the same property on all charts.

\begin{proof}
For one chart, take finitely many generators of the quotient.
By FCD9, lift each to a homogeneous element of $M$.  Let $N$
be their span.  It is a graded submodule: multiplying a homogeneous
generator by the homogeneous components of a polynomial decomposes
every element of $N$ into components that still belong to $N$.
Every $y\in M$ has $y-s\in(x-1)M$ for some $s\in N$.

First suppose $y$ has degree $n$.  Choose a common degree $D$
at least $n$ and all degrees occurring in $s$.  Homogenizing $s$
as in FCD9 gives $s_D\in N\cap M_D$.  The two elements
$x^{D-n}y$ and $s_D$ have the same image in $M/(x-1)M$.
Their difference is homogeneous, so FCD9 gives
$x^k(x^{D-n}y-s_D)=0$ for some $k$.  Hence
$x^{k+D-n}y=x^ks_D\in N$.  For a general element, apply
this argument to each of its finitely many homogeneous components
and increase to a common exponent.  Their sum then has the same
power-span property.

For finitely many charts take the union of the finite homogeneous
families just constructed.  In the more general formulation, first
take the union of the given finite families, and then replace each
element by all of its finitely many homogeneous components.  The
new span contains the old span on every chart.  The empty chart
index set is harmless: the empty family satisfies the vacuous
chart conditions.
\end{proof}
% END FCD10

% BEGIN FCD11
\paragraph{FCD11: the relations also have finite dehomogenizations.}
\label{claim:FCD11}
Let $\varphi:F\to G$ be a graded $S_R$-linear map with $F$
finitely generated.  Fix $i$, put $x=X_i$ and $u=x-1$, and
assume that every $z\in G$ has some $x^kz\in\operatorname{im}\varphi$.
If $G/uG$ is finitely presented, then
$(\ker\varphi)/u(\ker\varphi)$ is finitely generated.

\begin{proof}
Write $K=\ker\varphi$.  The composite
$F\to G\to G/uG$ is surjective: $x^kz$ and $z$ have the
same class modulo $uG$.  By FCD1 its kernel $J$ is finitely
generated.  We claim
\[
 J=K+uF,\qquad K\cap uF=uK.                          \tag{FCD.11}\label{FCD.11}
\]
The easy inclusion in the first equality follows by applying
$\varphi$.  For its converse, suppose $\varphi(y)=uw$.
By FCD9, write $y=y'+uv_0$ with $y'$ homogeneous.
Then $\varphi(y')=u(w-\varphi(v_0))$ is homogeneous.
FCD9 gives $x^k\varphi(y')=0$.  Therefore $k'=x^ky'$
belongs to $K$.  Write $x^ky'=y'+uv_1$ by the geometric-sum
identity; then $y=k'+u(v_0-v_1)$, as required.

For the second equality, if $uv\in K$, then
$u\varphi(v)=0$.  Multiplication by $u$ is injective on the
graded module $G$, by FCD9, so $v\in K$.  The reverse inclusion
is immediate.  The isomorphism theorem applied to \eqref{FCD.11}
now gives
\[
 K/uK\simeq (K+uF)/uF=J/uF.
\]
This is a quotient of the finitely generated module $J$ and is
therefore finitely generated.  No Noetherian hypothesis on $S_R$
or finite-presentation hypothesis on $F$ was used.
\end{proof}
% END FCD11

% BEGIN FCD12
\paragraph{FCD12: the full finite graded presentation up to chart torsion.}
\label{claim:FCD12}
Let $G$ be a graded $S_R$-module such that each
$G/(X_i-1)G$ is finitely presented over $S_R$.
There exist finite ranks $a,b$, integer degree lists $d,e$, a
matrix $\Phi$ with $\Phi_{lj}$ homogeneous of degree $e_j-d_l$,
and a graded map
\[
 \psi:N=\operatorname{coker}\Phi\longrightarrow G
\]
with the following two properties for every $i$:
\begin{align}
 \psi(y)=0&\ \Longrightarrow\ X_i^ky=0
       \quad\hbox{for some }k\ge0,\nonumber\\
 z\in G&\ \Longrightarrow\ X_i^kz\in\operatorname{im}\psi
       \quad\hbox{for some }k\ge0.                  \tag{FCD.12}\label{FCD.12}
\end{align}
These are the original quantifiers of
\texttt{exists\_\allowbreak{}graded\_\allowbreak{}presentation\_\allowbreak{}saturating}.

\begin{proof}
FCD10 gives finitely many homogeneous $g_l\in G_{d_l}$ whose
span has the second property in \eqref{FCD.12}.  Let
$F=\bigoplus_lS_Re_l$, with $\deg e_l=d_l$, and define the
graded map $\varphi(e_l)=g_l$.  Its kernel $K$ is graded:
if a sum of homogeneous components maps to zero, the direct-sum
decomposition in $G$ makes each component map to zero.

For every chart, FCD11 proves that $K/(X_i-1)K$ is finitely
generated.  Apply FCD10 to $K$ to choose a finite homogeneous
family $\kappa_j\in K_{e_j}$ whose span is correct up to powers
of every $X_i$.  Let $\Phi$ have these columns.  In the chosen
graded free basis its $l,j$ entry has degree $e_j-d_l$.
Since all columns lie in $K$, the map $\varphi$ descends to the
displayed graded $\psi:N\to G$.

Its image is exactly the image of $\varphi$, giving the second
assertion of \eqref{FCD.12}.  If $\psi(y)=0$, represent $y$
by $v\in F$.  Then $v\in K$.  For the chosen chart $i$, some
$X_i^kv$ is a combination of the $\kappa_j$, hence is in
$\Phi S_R^b$.  Its class $X_i^ky$ is zero.  This proves the
first assertion, for arbitrary $y$, not just homogeneous $y$.

This constructs the finite presentation needed on the charts without
asserting that $G$ itself is a finite $S_R$-module.
\end{proof}
% END FCD12

% BEGIN FCD13
\paragraph{FCD13: actual dehomogenization gives the required finite quotients.}
\label{claim:FCD13}
Let $G$ be graded, $\pi:S_R\twoheadrightarrow B$ an
$R$-algebra map with $\pi(X_i)=1$, and
$\epsilon:G\twoheadrightarrow P$ a $\pi$-semilinear map.
Suppose every homogeneous $y$ with $\epsilon(y)=0$ is killed
by some power of $X_i$.  Then
\[
 \ker\epsilon=(X_i-1)G,
 \qquad G/(X_i-1)G\simeq P                          \tag{FCD.13}\label{FCD.13}
\]
as $S_R$-modules, using $\pi$ for the action on $P$.
In particular the quotient is finite if $P$ is finite over $B$.
It is finitely presented over $S_R$ if $B$ is finitely presented
over $R$ and $P$ is finitely presented over $B$.

\begin{proof}
The inclusion $(X_i-1)G\subseteq\ker\epsilon$ follows from
semilinearity.  If $\epsilon(y)=0$, write
$y=y_N+(X_i-1)v$ with $y_N$ homogeneous, using FCD9.
Then $\epsilon(y_N)=0$, so $X_i^ky_N=0$ for some $k$.
The geometric-sum identity gives
\[
 y_N=-(X_i-1)(1+X_i+\cdots+X_i^{k-1})y_N,
\]
and $y$ also belongs to $(X_i-1)G$.  Surjectivity of
$\epsilon$ proves the isomorphism in \eqref{FCD.13}.

Because $\pi$ is surjective, a finite set of $B$-module generators
of $P$ also generates it over $S_R$.  For finite presentation,
FCD1 gives finitely many generators $g_k$ of $\ker\pi$.
Take a finite presentation $B^t\to B^s\to P\to0$ and lift its
relation columns to $S_R^s$.  These lifted columns together with
$g_ke_l$ generate the kernel of $S_R^s\to P$: reduce any
kernel vector to $B^s$, subtract the lifted relation combination,
and then express the remaining coordinates in $(g_k)$.
This gives an actual finite presentation over $S_R$ and proves
the assertion.
\end{proof}
% END FCD13

% BEGIN FCD14
\paragraph{FCD14: the same dehomogenization after the finite presentation.}
\label{claim:FCD14}
Let $\psi:N\to G$ be graded, and suppose it satisfies both
conditions \eqref{FCD.12} for a fixed $i$.  Let
$\epsilon:G\twoheadrightarrow P$ be semilinear for a ring map
$\pi:S_R\to B$ with $\pi(X_i)=1$, and suppose its homogeneous
kernel is killed by powers of $X_i$.  Then $\epsilon\psi$ is
surjective and
\[
 \epsilon\psi(y)=0\quad\Longleftrightarrow\quad
 y=(X_i-1)v\text{ for some }v\in N.                 \tag{FCD.14}\label{FCD.14}
\]
Only $N$ needs a direct-sum grading for this assertion; for $G$
it suffices that $\psi$ takes homogeneous elements to the stated
homogeneous pieces.

\begin{proof}
Given $p\in P$, choose $z\in G$ with $\epsilon(z)=p$.
By the second part of \eqref{FCD.12}, there are $k,y$ with
$\psi(y)=X_i^kz$.  Semilinearity and $\pi(X_i)=1$ give
$\epsilon\psi(y)=p$.

If $y$ is homogeneous and $\epsilon\psi(y)=0$, the chart
kernel hypothesis gives $X_i^k\psi(y)=0$.
Thus $\psi(X_i^ky)=0$, and the first part of
\eqref{FCD.12} gives $X_i^{j+k}y=0$ for some $j$.
The proof of FCD13's kernel identity, applied to the graded
module $N$ and to $\epsilon\psi$, now proves
\eqref{FCD.14}.  That proof only uses homogeneous pieces in
the target to invoke the stated kernel hypothesis, as claimed.
\end{proof}
% END FCD14

\subsection{The Noetherian graded-cube model}

For a graded $S_R$-module $M$ and a degree $n$, define its
localization cube by
\[
 C(M,n)_J=(M[X_J^{-1}])_n,\qquad J\subseteq\iota,    \tag{FCD.15}\label{FCD.15}
\]
with the actual localization maps for inclusions of subsets.  At
$J=\varnothing$ the entry is $M_n$.  This is the source's
\texttt{grCube}; the degree pieces are its \texttt{homLoc}.
The ordinary \v Cech complex has
\[
 \check C^p(M,n)=\bigoplus_{|J|=p+1}C(M,n)_J,
\]
and alternating localization differentials.  We also use the
augmented cube complex $A^p(M,n)=\bigoplus_{|J|=p}C(M,n)_J$,
whose degree-zero term is $M_n$.  These conventions will keep the
top and bottom cohomology in the next argument explicit.

% BEGIN FCD15
\paragraph{FCD15: finite graded \v Cech cohomology, with its monomial proof.}
\label{claim:FCD15}
Suppose $R_0$ is Noetherian, $\iota$ is finite, and $M$ is a
finite graded $S_{R_0}$-module.  For every integer $n$ and
every $p\ge0$, the cohomology $H^p(\check C(M,n))$ is a finite
$R_0$-module.

\begin{proof}
Put $m=|\iota|$.  First compute the augmented complex for one
free generator of degree $d$.  The localized degree-$n$ terms
have Laurent monomial bases $X^\alpha$ with
\[
 \alpha\in\mathbb Z^{\iota},\qquad
 \sum_i\alpha_i=n-d,
 \qquad\{i:\alpha_i<0\}\subseteq J.
\]
Indeed localizing a polynomial ring at $X_J$ allows precisely
the negative exponents in $J$.  Every element has a unique finite
such expansion, over arbitrary $R_0$, so the complex is the
direct sum of its complexes for individual exponent vectors.

Fix $\alpha$ and put $N(\alpha)=\{i:\alpha_i<0\}$.
Its summands are one copy of $R_0$ for every subset $J$ containing
$N(\alpha)$, in degree $|J|$.  With basis $e_J$, the differential
adds one index $j\notin J$ with sign
$(-1)^{\#\{i\in J:i<j\}}$.
If $N(\alpha)\ne\iota$, choose $j_0\notin N(\alpha)$.
Define a degree-minus-one map by
\[
 h(e_J)=
 \begin{cases}
 (-1)^{\#\{i\in J:i<j_0\}}e_{J\setminus\{j_0\}},&j_0\in J,\\
 0,&j_0\notin J.
 \end{cases}
\]
The indicated subset still contains $N(\alpha)$.
In $dh+hd$, the term removing and adding $j_0$ is $e_J$;
the terms involving any other index cancel in pairs, since
exchanging the order of insertion and deletion changes its sign.
Thus $dh+hd=1$ and this monomial complex is contractible.

If $N(\alpha)=\iota$, only the top term $J=\iota$ occurs,
giving one possible top cohomology class.  There are finitely many
such $\alpha$ for fixed $n-d$.  When $m>0$, writing
$\alpha_i=-1-\beta_i$ reduces the count to nonnegative integers
with fixed sum $\sum_i\beta_i=d-n-m$.  Each is bounded by that
sum when the sum is nonnegative, and there are none otherwise.
When $m=0$, there is one empty exponent vector if $n=d$ and
none otherwise.  Consequently every cohomology group of this
augmented free complex is finite.  Finite direct sums prove the
same assertion for finite graded free modules of arbitrary shifts.

Now let $M$ be any finite graded module.  It has a finite graded
free surjection $F\to M$: replace any finite generating family
by its homogeneous components.  Its kernel $K$ is graded and
finite because $S_{R_0}$ is Noetherian.  Localization is exact,
and taking a homogeneous component of an exact sequence of
graded modules is exact.  Hence
\[
 0\longrightarrow A^\bullet(K,n)
 \longrightarrow A^\bullet(F,n)
 \longrightarrow A^\bullet(M,n)\longrightarrow0
\]
is an exact sequence of complexes.  Prove finiteness of augmented
cohomology by descending induction on $p$, uniformly for all
finite graded modules.  Above $m$ the complexes are zero.  At
degree $p$, the long exact sequence gives
\[
 H^p(A(F,n))\longrightarrow H^p(A(M,n))
 \longrightarrow H^{p+1}(A(K,n)).
\]
The first module is finite by the free calculation; the last is
finite by the induction hypothesis.  The middle module is an
extension of a quotient of the first by a submodule of the last.
Both are finite because $R_0$ is Noetherian.  This completes
the induction without any finite-projective-dimension assumption.

Finally $M_n$ is finite over $R_0$: in the chosen free surjection,
$F_n$ is a finite sum of polynomial pieces of fixed nonnegative
degree, each with finitely many monomials, and it surjects onto
$M_n$.  For $p\ge1$, ordinary \v Cech cohomology is the
augmented cohomology in degree $p+1$.  In degree zero there is
the exact sequence
\[
 0\longrightarrow\operatorname{im}(M_n\to A^1(M,n))
 \longrightarrow H^0(\check C(M,n))
 \longrightarrow H^1(A(M,n))\longrightarrow0.
\]
Both outer modules are finite.  This proves the claim, including
the case $\iota=\varnothing$, when the ordinary complex is zero.
\end{proof}
% END FCD15

% BEGIN FCD16
\paragraph{FCD16: flatness on a singleton implies flatness on every containing vertex.}
\label{claim:FCD16}
For any graded $S_R$-module $M$ there is an $R$-linear
isomorphism
\[
 (M[X_i^{-1}])_0\simeq M/(X_i-1)M.                  \tag{FCD.16}\label{FCD.16}
\]
If this module is $R$-flat, then for every $J$ containing $i$
and every integer $n$, $(M[X_J^{-1}])_n$ is $R$-flat.

\begin{proof}
Represent a degree-zero element on the left of \eqref{FCD.16}
by $y/X_i^k$, with $k\ge0$ and $y\in M_k$, and send it to
the class of $y$.  Any degree-zero fraction has such a
representative: choose a common denominator and take its
homogeneous numerator of the required degree.  Equality of two
fractions means that an additional power of $X_i$ kills their
cross-multiplied difference.  Modulo $X_i-1$, this is exactly
equality of the two numerator classes, so the map is well-defined.

If the class of $y$ is zero, then $y=(X_i-1)w$ and is
homogeneous.  FCD9 gives $X_i^ly=0$, so the original fraction
was zero.  For surjectivity, take any class represented by $z$;
FCD9 gives a homogeneous representative $z_N\in M_N$ for
some $N\ge0$.  The fraction $z_N/X_i^N$ maps to it.

Multiplication by the invertible homogeneous element $X_i$
identifies consecutive degrees of $M[X_i^{-1}]$.
Therefore flatness of its degree-zero piece gives flatness of
every piece, and their direct sum $M[X_i^{-1}]$ is $R$-flat.
Further localize by the product of the remaining $X_j$ in $J$.
Localization of an $R$-flat module by an endomorphism commuting
with $R$ is $R$-flat: it is the filtered direct limit of copies
of that module with transition given by that endomorphism, and
filtered direct limits preserve injections after tensoring.
This proves flatness of $M[X_J^{-1}]$.
Its grading is a direct sum, so every homogeneous piece is a
direct summand and is flat as well.  The grading and its directness
can be checked by multiplying a finite sum of homogeneous fractions
by a common homogeneous denominator and using directness in $M$.
\end{proof}
% END FCD16

% BEGIN FCD17
\paragraph{FCD17: canonical scalar extension and localization comparisons.}
\label{claim:FCD17}
Let $R_0\to R$ be any ring map and let a homogeneous matrix
$\Phi_0$ over $S_{R_0}$ map to $\Phi$ over $S_R$, with the
same free-module shifts.  Write $N_0=\operatorname{coker}\Phi_0$
and $N=\operatorname{coker}\Phi$.  There is a canonical
isomorphism of cubes, for every degree $n$,
\[
 R\otimes_{R_0} C(N_0,n)\simeq C(N,n),               \tag{FCD.17}\label{FCD.17}
\]
induced by coefficientwise application of the ring map to the
representatives.  If $\psi:N\to G$ is graded and satisfies
\eqref{FCD.12} for every $i$, its induced map
$C(N,n)_J\to C(G,n)_J$ is bijective for every nonempty $J$.
All these maps commute with the restriction maps between vertices.

\begin{proof}
The map $R\otimes_{R_0}S_{R_0}\to S_R$ sending
$r\otimes p$ to $r\,p_R$ is an isomorphism: the monomials
are a basis on both sides.  Apply right exactness of tensor
product to the presentation
$S_{R_0}^b\xrightarrow{\Phi_0}S_{R_0}^a\to N_0\to0$.
It gives the canonical isomorphism
$R\otimes_{R_0}N_0\simeq N$, sending
$1\otimes[v]$ to $[v_R]$.  This uses a cokernel, not a
kernel, so flatness of $R$ over $R_0$ is unnecessary.

Localizing commutes with this scalar extension.  Explicitly the
map sends $r\otimes(y/X_J^k)$ to $r\,y_R/X_J^k$;
the inverse follows from the universal property of inverting
$X_J$, or equivalently from the same filtered-limit description
used in FCD16.  These maps preserve degrees.  Taking a degree
piece also commutes with tensor product here: a graded module
is the direct sum of its pieces, tensor product commutes with
that direct sum, and its coordinate projections split the
individual pieces.  This proves \eqref{FCD.17}, including the
empty vertex.

For the second assertion, choose $i\in J$.  The kernel and
cokernel of $\psi$ have every element killed by a power of
$X_i$, by \eqref{FCD.12}.  Localizing at $X_J$ kills both,
because $X_i$ becomes invertible.  Exactness of localization
then gives an isomorphism $N[X_J^{-1}]\simeq G[X_J^{-1}]$,
and its degree-$n$ component is bijective.
Each map sends a numerator to its specified image and leaves
the localization denominators in place.  Enlarging $J$ has
the same effect before or after applying these formulas.
This proves the asserted commutation and hence the cube,
rather than merely vertexwise, statement.
\end{proof}
% END FCD17

% BEGIN FCD18
\paragraph{FCD18: the original Noetherian graded-cube model.}
\label{claim:FCD18}
Let $R$ be arbitrary and $G$ a graded $S_R$-module.  For each
$i\in\iota$, suppose there exist a smooth $R$-algebra $B_i$,
a surjective $R$-algebra map $\pi_i:S_R\to B_i$ with
$\pi_i(X_i)=1$, a finitely presented projective $B_i$-module
$P_i$, and a surjective $\pi_i$-semilinear map
$\epsilon_i:G\to P_i$.  Assume that whenever $y$ is
homogeneous and $\epsilon_i(y)=0$, some $X_i^ky$ is zero.
Then there are a single Noetherian subring $R_0\subseteq R$,
a cube $C_0$ of $R_0$-modules, and an actual cube map
\[
 \theta:R\otimes_{R_0}C_0\longrightarrow C(G,0)       \tag{FCD.18}\label{FCD.18}
\]
such that:
\begin{enumerate}
\item $(C_0)_J$ is $R_0$-flat for every nonempty $J$;
\item every $H^p(\check C(C_0))$ is finite over $R_0$;
\item $\theta_J$ is bijective for every nonempty $J$.
\end{enumerate}
This is the full original
\texttt{exists\_\allowbreak{}noetherian\_\allowbreak{}model\_\allowbreak{}gr\allowbreak{}Cube}, with its original
chart hypotheses and quantifier order.

\begin{proof}
FCD13 makes every $G/(X_i-1)G$ finitely presented over $S_R$.
FCD12 therefore supplies one homogeneous matrix $\Phi$ and
one graded $\psi:N=\operatorname{coker}\Phi\to G$ satisfying
\eqref{FCD.12} on all charts.
Let $q:S_R^a\to N$ be its quotient map.
For each $i$, FCD14 says that $\epsilon_i\psi$ is surjective
with kernel $(X_i-1)N$.  Its composite with $q$ consequently
has the exact kernel
\[
 \ker(\epsilon_i\psi q)=\Phi S_R^b+(X_i-1)S_R^a.
\]
Indeed $q(y)=(X_i-1)q(v)$ is equivalent to
$y-(X_i-1)v\in\Phi S_R^b$.  Thus the ORIGINAL $B_i,P_i$
and these composite maps satisfy every hypothesis of FCD8.

Apply FCD8 to obtain one $R_0$ and $\Phi_0$.  Give
$N_0=\operatorname{coker}\Phi_0$ the same quotient grading,
and define, rather than choose independently,
\[
 C_0=C(N_0,0).
\]
FCD17 supplies the canonical first map in the composite
\[
 R\otimes_{R_0}C(N_0,0)
 \xrightarrow{\ \simeq\ } C(N,0)
 \xrightarrow{C(\psi,0)} C(G,0).
\]
Define $\theta$ to be exactly this composite.

For a nonempty $J$, choose $i\in J$.
FCD8 makes $N_0/(X_i-1)N_0$ flat over $R_0$; FCD16
therefore makes $(N_0[X_J^{-1}])_0=(C_0)_J$ flat.
The module $N_0$ is finite over the Noetherian polynomial ring
$S_{R_0}$, so FCD15 gives finite \v Cech cohomology in every
degree.  Finally FCD17 makes both factors of $\theta_J$
bijective for every nonempty $J$.  The maps commute with all
restrictions by that same claim, so the conclusion concerns
the actual cube map, as required.

The choices $\Phi,\psi,R_0,\Phi_0,C_0,\theta$ have all
been made for the fixed original $G$ and its chart data.
No later coefficient algebra has been introduced.  In particular,
the original map $R_0\to R$ need not be flat.
\end{proof}
% END FCD18

\subsection{Exact use in the existing section proof}

FCD18 now supplies the Noetherian-cube input previously retained in
GSL8 whenever its original graded chart hypotheses are available.
Applying GSL8 to $C(G,0)$ gives one bounded complex of finite
projective $R$-modules and one map to its \v Cech complex that
remains a quasi-isomorphism after every $R\to R'$.  This is a
direct reuse of the already proved complex algebra.

For the geometric application, the intended $G$ is the graded sum of
the actual global sections of $M\otimes H^t$, and the $X_i$ act by
the fixed sections $s_i$ of the very ample bundle $H$.
The remaining producers must identify the chart algebra with
$\Gamma(D(s_i),\mathcal O)$ and its module with the actual
sections of $M\otimes H^0$ there, prove the required smoothness,
finite presentation and projectivity, and prove the dehomogenization
and actual-section comparison.  They must produce the maps on this
same graded module and the same cover.  After these are supplied,
the source's final map is the composite of $\theta$, the geometric
localization comparison, and the natural identification
$M\otimes H^0=M$; GSL8--GSL10 then give every literal
\texttt{Perf} field.  The current section does not replace these
geometric arguments by an abstract identification of dimensions.

The other retained GSL producers---polarized base change, geometric
$H^1$ vanishing, geometric integrality and principal rank one,
torsion-section decomposition and growth, and the multiplication
and square generators---are unchanged.  The original \texttt{Moduli}
record and the HUF/PUN foundations also remain.  Neither a global
section law nor a selected whole RH/GRH obligation is closed by the
present algebraic model theorem.

\subsection{Source correspondence}

All following paths are relative to
\texttt{src/Development/AbelianVarietyTheory/Proofs/} at the
accepted source pin.  In this list \texttt{C/} abbreviates
\texttt{CoherentCohomologyProofs/}.  The local snapshots retain the
complete files; the provenance records bind the stated ranges and
each numbered proof block.

\begin{itemize}
\item FCD1--FCD3 and FCD7--FCD8:
\texttt{C/\allowbreak{}Exists\allowbreak{}Noetherian\allowbreak{}Subring\allowbreak{}Homogeneous\allowbreak{}Matrix\allowbreak{}Flat\allowbreak{}Char.lean},
lines 11--142 and 186--287.  FCD1 expands the finite-presentation
facts used at the start of the chart-witness proof.
\item FCD4:
\texttt{C/\allowbreak{}Exists\allowbreak{}Witnesses\allowbreak{}Of\allowbreak{}Formally\allowbreak{}Smooth\allowbreak{}Quotient.lean},
lines 9--72, and the converse certificate in
\texttt{C/\allowbreak{}Exists\allowbreak{}Noetherian\allowbreak{}Subring\allowbreak{}Homogeneous\allowbreak{}Matrix\allowbreak{}Flat\allowbreak{}Char.lean},
lines 145--182.
\item FCD5 expands the Noetherian flatness consequence used in that
same file, lines 253--268.  The dependency pin recorded by the
repository's \texttt{lake-manifest.json} is Mathlib
\texttt{0df444a360eaa60ab8c11dca51a86af692955474}; the supporting
files are \texttt{Mathlib/\allowbreak{}Ring\allowbreak{}Theory/\allowbreak{}Smooth/\allowbreak{}Flat.lean},
\texttt{Smooth/\allowbreak{}Adic\allowbreak{}Completion.lean}, and
\texttt{Adic\allowbreak{}Completion/\allowbreak{}As\allowbreak{}Tensor\allowbreak{}Product.lean} under
\texttt{Mathlib/RingTheory/}.
\item FCD6:
\texttt{C/\allowbreak{}Flat\allowbreak{}Quot\allowbreak{}SMul\allowbreak{}Top\allowbreak{}Coker\allowbreak{}Of\allowbreak{}Matrix\allowbreak{}Witnesses.lean}, lines 10--117.
\item FCD9--FCD12:
\texttt{C/\allowbreak{}Exists\allowbreak{}XPow\allowbreak{}Smul\allowbreak{}Eq\allowbreak{}Zero\allowbreak{}Of\allowbreak{}Eq\allowbreak{}XSub\allowbreak{}One\allowbreak{}Smul.lean},
\texttt{C/\allowbreak{}Exists\allowbreak{}XPow\allowbreak{}Smul\allowbreak{}Mem\allowbreak{}Span\allowbreak{}Of\allowbreak{}Finite\allowbreak{}Quot.lean},
\texttt{C/\allowbreak{}Exists\allowbreak{}Homogeneous\allowbreak{}Family\allowbreak{}Saturating.lean},
\texttt{C/\allowbreak{}Exists\allowbreak{}Pow\allowbreak{}Smul\allowbreak{}Eq\allowbreak{}Add\allowbreak{}Sub\allowbreak{}One\allowbreak{}Smul.lean}, and
\texttt{C/\allowbreak{}Exists\allowbreak{}Graded\allowbreak{}Presentation\allowbreak{}Saturating.lean}.
These are respectively the degree recurrence, finite chart generation,
homogeneous union, finite relation module and full presentation.
\item FCD13--FCD14:
\texttt{C/\allowbreak{}Dehomogenising\allowbreak{}Eq\allowbreak{}Zero\allowbreak{}Iff.lean},
\texttt{C/\allowbreak{}Finite\allowbreak{}Quot\allowbreak{}SMul\allowbreak{}Top\allowbreak{}XSub\allowbreak{}One\allowbreak{}Of\allowbreak{}Chart.lean}, and
\texttt{C/\allowbreak{}Dehomogenising\allowbreak{}Comp\allowbreak{}Of\allowbreak{}Saturating.lean}, lines 22--42.
\item FCD15:
\texttt{C/\allowbreak{}Cube\allowbreak{}Complex\allowbreak{}Poly\allowbreak{}Grading\allowbreak{}Finite\allowbreak{}Exact.lean},
\texttt{C/\allowbreak{}Cube\allowbreak{}Complex\allowbreak{}Gr\allowbreak{}Free\allowbreak{}Finite\allowbreak{}And\allowbreak{}Exact\allowbreak{}At.lean},
\texttt{C/\allowbreak{}Finite\allowbreak{}Homology\allowbreak{}Cube\allowbreak{}Complex\allowbreak{}Graded\allowbreak{}Module.lean},
\texttt{C/\allowbreak{}Finite\allowbreak{}Gr\allowbreak{}Cech\allowbreak{}Homology.lean}, and the explicit monomial
and \v Cech comparison declarations in the manifest.
\item FCD16--FCD17:
\texttt{C/\allowbreak{}Nonempty\allowbreak{}Hom\allowbreak{}Loc\allowbreak{}Singleton\allowbreak{}Equiv\allowbreak{}Quot\allowbreak{}SMul\allowbreak{}Top.lean},
\texttt{C/\allowbreak{}Flat\allowbreak{}Hom\allowbreak{}Loc\allowbreak{}Of\allowbreak{}Flat\allowbreak{}Hom\allowbreak{}Loc\allowbreak{}Singleton.lean},
\texttt{C/\allowbreak{}Is\allowbreak{}Base\allowbreak{}Change\allowbreak{}Coker\allowbreak{}Matrix\allowbreak{}Mv\allowbreak{}Polynomial.lean}, and
\texttt{C/\allowbreak{}Nonempty\allowbreak{}Base\allowbreak{}Change\allowbreak{}Gr\allowbreak{}Cube\allowbreak{}Iso.lean}.
\item FCD18:
\texttt{C/\allowbreak{}Dehomogenising\allowbreak{}Comp\allowbreak{}Of\allowbreak{}Saturating.lean}, lines 44--114.
The next geometric consumer is
\texttt{C/\allowbreak{}Exists\allowbreak{}Finite\allowbreak{}Projective\allowbreak{}Cech\allowbreak{}Complex\allowbreak{}Line\allowbreak{}Bundle.lean},
lines 19--43; the already reconstructed algebra it invokes is
GSL8--GSL10.
\end{itemize}
